\section{From the state count to an isolated mass}
\label{sec:mass-support}

Each mass in the support of $\rho$ contributes to the continuum two-point
function throughout an interval of spatial momenta. We will test these
contributions by positive measures from the finite-circumference transfer
operator. Distinct masses will then require distinct low-energy states
at every allowed momentum in that interval. Proposition~\ref{cou:linear}
bounds how many such masses there can be.

\subsection{Measures at fixed momentum}

Fix nonzero nonnegative smooth functions $g,h$ of compact support, with
$\supp g\subset(0,l)$ for a fixed positive integer $l$, and set
$f(t,b)=g(t)h(b)$. We use color $3$ throughout. On a cylinder of
circumference $w=w_k$, let $X_f$ be the source-normalized lattice field
sum and define its bounded truncation by
\[
 F=\max\{-w,\min\{X_f,w\}\}.
\]
Use the slab $[0,l]$ for the normalized insertion $A_F$, with the initial
row as its left kernel variable, and put
\[
 v=A_F\Omega_{N,w}.
\]
With this convention, the matrix element
$\langle v,T_{N,w}^{n/a}v\rangle$ places the first copy of $F$ on
$[-l,0]$ by time reflection and the second copy on $[n,n+l]$ by time
translation. To see the corresponding time weights, let $Q_{h,N}$
denote row multiplication by
$B_Na\sum_{b\in a\Z/w\Z}h(b)q_b^3$, using the periodization of $h$.
Before truncation the linear insertion satisfies the exact formula
\[
 A_{X_f}
   =a\sum_{t\in a\Z\cap[0,l]}
      g(t)T^{t/a}Q_{h,N}T^{(l-t)/a},
 \qquad
 A_{X_f}\Omega
   =a\sum_t g(t)T^{t/a}Q_{h,N}\Omega.
\]

Work in the complex transfer Hilbert space. Spatial row translations
$U_b$, for $b\in a\Z/w\Z$, commute with $T_{N,w}$. Choose their character
convention so that the momentum projection is
\[
 \Pi_p=\frac{a}{w}\sum_{b\in a\Z/w\Z}e^{-ipb}U_b,
 \qquad
 p\in\frac{2\pi}{w}\Z\cap[-\pi/a,\pi/a).
\]
For each such $p$, define the finite positive measure
\begin{equation}\label{mass:finite-measure}
 \nu_{N,w,p}
   =w\sum_E\bigl\|P_E\Pi_pv\bigr\|^2\delta_{e^{-E}}
 \quad\text{on }[0,1],
\end{equation}
where $P_E$ is the transfer-energy projection. The factor $w$ converts
the normalized character sum into a spatial Riemann sum.

Let $C_f(n,b)$ be the continuum correlation of $\phi^3(f)$ reflected in
time with a second copy translated by $n$ in time and by $b$ in space.
For a compact momentum interval $I$, we will write $\max_{p\in I}$ for
the maximum over allowed lattice momenta lying in $I$.

\begin{lemma}\label{mass:moments}
For every integer $n\ge0$ and every compact interval $I$ with nonempty
interior,
\begin{equation}\label{mass:moments-eq}
 \lim_{\substack{w\to\infty\\w=w_k}}\limsup_{N\to\infty}
 \max_{p\in I}
 \left|
   \int x^n\,\nu_{N,w,p}(dx)
   -\int_\R e^{-ipb}C_f(n,b)\,db
 \right|=0.
\end{equation}
\end{lemma}

\begin{proof}
The spectral theorem and the character formula give
\begin{align*}
 \int x^n\,\nu_{N,w,p}(dx)
 &=w\langle v,T_{N,w}^{n/a}\Pi_pv\rangle\\
 &=a\sum_{\substack{b\in a\Z\\-w/2\le b<w/2}}
       e^{-ipb}\langle v,T_{N,w}^{n/a}U_bv\rangle.
\end{align*}
The matrix element in the sum is the cylinder correlation of the two
bounded slab observables just described. For fixed $f,n$ and large $w$,
their joint support has time extent at most $n+2l$ and spatial extent at
most $w/2+\operatorname{diam}(\supp h)$. Both are below $3w/4$.
The cylinder-to-torus comparison, followed by
Proposition~\ref{tra:comparison}, replaces this correlation by its cutoff
plane counterpart with error at most $Cw^2e^{-cw}$, uniformly in $b$ and
$N$. The uniform field tails then remove the two truncations on the
plane with error $Ce^{-c'w}$, by Cauchy--Schwarz and the moment majorants.
The factor $a$ times the number of spatial translates is $w$, so the
total error in the displayed Riemann sum still tends to zero exponentially
in $w$.

For a fixed $w$, the cutoff plane smeared correlations converge uniformly
for $|b|\le w/2$ to $C_f(n,b)$. To see the uniformity, moment convergence
gives pointwise convergence, while the difference-test bound
\eqref{inp:test-difference} gives uniform equicontinuity. More explicitly,
on a fixed compact translation interval the smooth translates satisfy
$|f_b-f_{b'}|\le C|b-b'|$ on a fixed finite union of unit boxes. The
second-moment majorant bounds the $L^2$ norm of the corresponding field
difference by $C'|b-b'|$. Cauchy--Schwarz then gives the same modulus of
continuity for the correlations. Compactness upgrades pointwise
convergence to uniform convergence. Consequently, as $N\to\infty$, the
Riemann sum of the cutoff plane correlations converges to the continuum
integral over $[-w/2,w/2]$, uniformly for $p$ in $I$; the exponential
factors have a uniform modulus of continuity on this compact set.

Finally, the separated-time kernel bound in Section~\ref{sec:inputs}
implies $|C_f(n,b)|\le Ce^{-c|b|}$ outside a fixed compact interval.
The omitted spatial tail therefore tends to zero uniformly in $p$ as
$w\to\infty$. This proves \eqref{mass:moments-eq}. The case $n=0$ is
included: $\supp g$ has positive distance from zero, so the two reflected
field supports remain separated in time.
\end{proof}

The Euclidean spectral representation \eqref{inp:euclidean} identifies the
limiting measures explicitly. Write
\[
 \widehat h(p)=\int_\R e^{-ipb}h(b)\,db,
 \qquad G(E)=\int_0^l g(t)e^{-tE}\,dt,
 \qquad E(p,\mu)=\sqrt{p^2+\mu^2}.
\]
For each $p\in\R$, let $\nu_p$ be the image, under
$\mu\mapsto e^{-E(p,\mu)}$, of the positive measure
\begin{equation}\label{mass:measure}
 \frac{|\widehat h(p)|^2}{2E(p,\mu)}
       |G(E(p,\mu))|^2\,\rho(d\mu^2).
\end{equation}
Integrating the kernel against the two slab tests and then Fourier
transforming in their relative spatial displacement gives
\[
 \int_\R e^{-ipb}C_f(n,b)\,db=\int x^n\,\nu_p(dx).
\]
The time separation supplies exponential damping in $\mu$, so the
integrals converge and depend continuously on $p$ on compact intervals.
This also justifies the Fourier calculation, either directly or first
as an identity of distributions and then of continuous functions.

We will need continuous spectral tests rather than moments. For every
$\Psi\in C(I\times[0,1])$, Lemma~\ref{mass:moments} implies
\begin{equation}\label{mass:continuous-tests}
 \lim_{\substack{w\to\infty\\w=w_k}}\limsup_{N\to\infty}
 \max_{p\in I}
 \left|
  \int\Psi(p,x)\,\nu_{N,w,p}(dx)
   -\int\Psi(p,x)\,\nu_p(dx)
 \right|=0.
\end{equation}
Indeed, the $n=0$ case uniformly bounds the total masses in this iterated
limit. Approximate $\Psi$ uniformly by polynomials in $x$ with continuous
coefficients in $p$, for instance its Bernstein polynomials in $x$.
Each coefficient is bounded on $I$, so Lemma~\ref{mass:moments} applies
to the finite sum; the uniform approximation controls the remainder
against both measures. Thus no mass near $x=0$ is lost in passing from
moments to continuous tests.

\subsection{Finiteness of the low mass support}

\begin{proposition}\label{mass:finite-support}
The set
\[
 \{\mu\in[m_*,5m_*/4]:\mu^2\in\supp\rho\}
\]
is finite.
\end{proposition}

\begin{proof}
Choose $p_0>0$ so small that
\[
 \sqrt{p_0^2+(5m_*/4)^2}<u_0,
 \qquad \widehat h(p)\ne0\quad (|p|\le p_0).
\]
This is possible because $5m_*/4<u_0=3m_*/2$ and
$\widehat h(0)=\int h>0$. This momentum interval is fixed before any
finite collection of masses is selected.

Choose $J$ distinct masses $\mu_1,\ldots,\mu_J$ from the indicated
support. On $I=[-p_0,p_0]$, the curves
\[
 x_j(p)=e^{-E(p,\mu_j)},\qquad j=1,\ldots,J,
\]
are pairwise disjoint compact graphs in $I\times(e^{-u_0},1)$. Their
mutual distances and their distances from the two horizontal boundary
lines are positive. There are therefore nonnegative continuous functions
$\Psi_j$ on $I\times[0,1]$, supported in mutually disjoint neighborhoods
of these graphs contained in $I\times(e^{-u_0},1)$, and equal to one on
the respective graphs.

For every $p\in I$,
\[
 L_j(p):=\int\Psi_j(p,x)\,\nu_p(dx)>0.
\]
Indeed, membership of $\mu_j^2$ in $\supp\rho$ gives positive measure to
every neighborhood of that point. The integrand in \eqref{mass:measure}
is positive there: $\widehat h(p)\ne0$ and $G(E)>0$ for every $E>0$,
since $g\ge0$ is nonzero. Continuity of $\Psi_j$ makes its value positive
through a neighborhood of $\mu_j$ at the fixed momentum. Moreover
$L_j$ is continuous in $p$, by exponential domination in
\eqref{mass:measure}. Compactness yields
\[
 \eta:=\min_{1\le j\le J}\min_{p\in I}L_j(p)>0.
\]
Both $\eta$ and the required approximation scale may depend on this
finite collection of masses.

By \eqref{mass:continuous-tests}, for all sufficiently large fixed
dyadic $w$, and then all sufficiently large $N$ depending on $w$,
\[
 \int\Psi_j(p,x)\,\nu_{N,w,p}(dx)>\eta/2
 \qquad (j=1,\ldots,J)
\]
at every allowed $p\in I$. Each atom of the measure
\eqref{mass:finite-measure} comes from a transfer-energy eigenspace
with momentum $p$. The disjoint
supports of the $\Psi_j$ therefore require at least $J$ distinct energies
below $u_0$ in that space. Distinct momentum spaces are orthogonal.
For large $w$, the number of momenta in $I$ is at least
$c_0w$, where $c_0>0$ depends only on $p_0$; at any fixed $w$ these are
distinct Brillouin characters once $N$ is large. Thus
\[
 Jc_0w\le\mathcal N_{N,w}(u_0).
\]
Taking the cutoff limit and using Proposition~\ref{cou:linear} gives
$Jc_0\le C$. The constants $c_0,C$ are independent of the selected
masses and of $J$. The support can consequently contain only finitely
many points in the stated interval.
\end{proof}

\begin{proof}[Proof of Theorem~\ref{intro:main}]
The nonempty closed mass support has minimum $m_*>0$, by the continuum
inputs in Section~\ref{sec:inputs}. Proposition~\ref{mass:finite-support}
makes this minimum isolated in the support. A positive locally finite
measure assigns strictly positive finite mass to an isolated support
point, so
\[
 Z=\rho(\{m_*^2\})\in(0,\infty).
\]
Choose $\delta>0$ small enough that no other mass in the support lies
below $m_*+\delta$, and define
\[
 m_1=m_*,\qquad \rho_{\mathrm{rest}}=\rho-Z\delta_{m_*^2}.
\]
This is a positive measure, and its support is contained in
$[(m_1+\delta)^2,\infty)$, as required.
\end{proof}
